% !TEX root = ../main.tex
\corrigesChapitre{16}{Mouvement dans un champ de gravitation}
\corrigesSouples

\begin{blocLivre}{cor-16-1}
\begin{corrige}{1}{La Terre et la Lune}
\demarche{Loi de la gravitation, puis deuxième loi de Newton appliquée à la Lune.}
\begin{questions}
\item $F=G\,\dfrac{M_\mathrm{T}M_\mathrm{L}}{d^2}=\res{\qty{1,98e20}{N}}$ ; la Lune exerce sur la Terre une force
\res{\text{de même valeur}}, de sens opposé.
\item $a=\dfrac{F}{M_\mathrm{L}}=\res{\qty{2,70e-3}{m.s^{-2}}}$, dirigée vers le centre de la Terre.
\item $\mathcal{G}_\mathrm{L}=G\,\dfrac{M_\mathrm{L}}{R_\mathrm{L}^2}=\res{\qty{1,62}{N.kg^{-1}}}$, environ \res{\text{six fois moins}} que sur Terre.
\end{questions}
\end{corrige}
\end{blocLivre}
\begin{blocLivre}{cor-16-2}
\begin{corrige}{2}{Un satellite en orbite basse}
\demarche{Référentiel géocentrique ; deuxième loi dans le repère de Frenet ; $r=R_\mathrm{T}+h$.}
\begin{questions}
\item Dans le référentiel \res{\text{géocentrique}}, supposé galiléen ; seule la force de gravitation exercée par la Terre
s'exerce sur le satellite.
\item $m\,\vec{a}=G\,\dfrac{m\,M_\mathrm{T}}{r^2}\,\un$ : selon $\ut$, $\ddt{v}=0$ (uniforme) ; selon $\un$,
$\dfrac{v^2}{r}=\dfrac{GM_\mathrm{T}}{r^2}$, d'où $\res{v=\sqrt{GM_\mathrm{T}/r}}$.
\item $r=\qty{6921}{km}$ : $v=\res{\qty{7,58e3}{m.s^{-1}}}$ et $T=\dfrac{2\pi r}{v}=\res{\qty{5,73e3}{s}}$, soit 95,6~min.
\item \res{\text{Oui}} : $v$ et $T$ ne dépendent pas de la masse du satellite.
\end{questions}
\end{corrige}
\end{blocLivre}
\begin{blocLivre}{cor-16-3}
\begin{corrige}{3}{L'orbite d'une comète}
\demarche{Le Soleil est à un foyer ; le rayon balaie des aires égales en des durées égales.}
\begin{questions}
\item L'orbite est une ellipse dont le Soleil occupe \res{\text{l'un des foyers}}, ici le foyer le plus proche de P.
\item Vitesse \res{\text{maximale au périhélie P}}, le plus proche du Soleil, \res{\text{minimale à l'aphélie A}} : pour
balayer la même aire, la comète parcourt un arc plus long quand elle est proche.
\item Les durées étant égales, \res{\text{les aires sont égales}} ; les arcs parcourus sont pourtant très différents.
\item Le Soleil est au centre, et la loi des aires impose une vitesse constante : mouvement \res{\text{circulaire uniforme}} ;
la troisième loi s'écrit avec le rayon $r$.
\end{questions}
\end{corrige}
\end{blocLivre}
\begin{blocLivre}{cor-16-4}
\begin{corrige}{4}{Le satellite géostationnaire}
\demarche{$T$ égale à la rotation de la Terre dans le référentiel géocentrique, puis troisième loi.}
\begin{questions}
\item Orbite circulaire \res{\text{dans le plan de l'équateur}}, parcourue \res{\text{dans le sens de rotation de la Terre}},
avec une période \res{\text{égale à la durée de rotation de la Terre}}.
\item \qty{86400}{s} est la durée du jour solaire, liée au Soleil ; dans le référentiel géocentrique, la Terre fait un tour en
\qty{86164}{s} (jour sidéral).
\item $r=\left(\dfrac{GM_\mathrm{T}T^2}{4\pi^2}\right)^{1/3}=\res{\qty{4,21e7}{m}}$, soit une altitude de \res{\num{35800}~\text{km}}.
\item $v=\dfrac{2\pi r}{T}=\res{\qty{3,07e3}{m.s^{-1}}}$. Le satellite reste à la verticale du même point : vu du sol, il est
immobile.
\end{questions}
\end{corrige}
\end{blocLivre}
\begin{blocLivre}{cor-16-5}
\begin{corrige}{5}{Peser Jupiter}
\demarche{Troisième loi : $\dfrac{T^2}{r^3}=\dfrac{4\pi^2}{GM}$, donc $M=\dfrac{4\pi^2r^3}{GT^2}$.}
\begin{questions}
\item Dans un référentiel lié au centre de Jupiter (« jovicentrique »), supposé galiléen.
\item $\dfrac{T^2}{r^3}=\dfrac{4\pi^2}{G\,M_\mathrm{J}}$.
\item Avec $T=1{,}769\times86400=\qty{1,528e5}{s}$ : $M_\mathrm{J}=\res{\qty{1,90e27}{kg}}$, soit \res{\text{319 fois}} la masse de la Terre.
\end{questions}
\end{corrige}
\end{blocLivre}
\begin{blocLivre}{cor-16-6}
\begin{corrige}{6}{La comète de Halley}
\demarche{Aux extrémités du grand axe, l'aire balayée pendant $\Delta t$ est un triangle de base $r$ et de hauteur $v\,\Delta t$.}
\begin{questions}
\item Pendant $\Delta t$, l'aire balayée vaut $\frac12\,r\,v\,\Delta t$ ; la loi des aires donne
$\frac12\,r_\mathrm{P}v_\mathrm{P}\Delta t=\frac12\,r_\mathrm{A}v_\mathrm{A}\Delta t$, d'où l'égalité.
\item $v_\mathrm{A}=v_\mathrm{P}\,\dfrac{r_\mathrm{P}}{r_\mathrm{A}}=54{,}5\times\dfrac{0{,}586}{35{,}1}=\res{\qty{0,91}{km.s^{-1}}}$.
\item $a=\dfrac{0{,}586+35{,}1}{2}=\res{\qty{17,8}{ua}}$.
\item $T^2=a^3$ (en années et ua) : $T=17{,}84^{3/2}=\res{75{,}4~\text{ans}}$ ; retour prévu en \res{2061}.
\end{questions}
\end{corrige}
\end{blocLivre}
\begin{blocLivre}{cor-16-7}
\begin{corrige}{7}{Tester la troisième loi en Python}
\demarche{Le rapport $T^2/a^3$ est constant si la troisième loi est vérifiée.}
\begin{questions}
\item Un \res{\text{dictionnaire}} : à chaque nom de planète est associé un couple (demi-grand axe en ua, période en années).
\item Pour la Terre, $a=\qty{1}{ua}$ et $T=\qty{1}{an}$ par définition des unités. Pour les autres, le rapport vaut 1 aux arrondis
des données près : \res{\text{la troisième loi est vérifiée}}.
\item On construit les listes \code{A = [a**3 for (a, T) in planetes.values()]} et
\code{B = [T**2 for (a, T) in planetes.values()]}, puis \code{plt.plot(A, B, "o")} et \code{plt.show()} : les points
sont alignés sur une \res{\text{droite passant par l'origine}}.
\item $\dfrac{T^2}{a^3}=\dfrac{(\num{3.156e7})^2}{(\num{1.496e11})^3}=\dfrac{4\pi^2}{GM_\odot}$, d'où $M_\odot=\res{\qty{1,99e30}{kg}}$.
\end{questions}
\schema[corrige]{rev-16-ex07-cor}
\legende{Les six planètes du programme sont compatibles avec une droite passant par l’origine.}
\end{corrige}
\end{blocLivre}
\begin{blocLivre}{cor-16-8}
\begin{corrige}{8}{Peser la Terre grâce à la Lune}
\demarche{Même démarche que pour un satellite artificiel, avec $T=27{,}32$ jours.}
\begin{questions}
\item Deuxième loi dans le repère de Frenet : $\dfrac{v^2}{r}=\dfrac{GM_\mathrm{T}}{r^2}$, d'où $v=\sqrt{\dfrac{GM_\mathrm{T}}{r}}$.
\item $T=\dfrac{2\pi r}{v}=2\pi\sqrt{\dfrac{r^3}{GM_\mathrm{T}}}$, d'où $\dfrac{T^2}{r^3}=\dfrac{4\pi^2}{GM_\mathrm{T}}$.
\item $M_\mathrm{T}=\dfrac{4\pi^2r^3}{GT^2}=\res{\qty{6,03e24}{kg}}$.
\item Écart de 1~\% : la Lune n'est pas négligeable devant la Terre (1/81 de sa masse), et son orbite n'est pas exactement
circulaire.
\end{questions}
\end{corrige}
\end{blocLivre}
\begin{blocLivre}{cor-16-9}
\begin{corrige}{9}{Les satellites GPS}
\demarche{$r=R_\mathrm{T}+h$, puis $T=2\pi\sqrt{r^3/(GM_\mathrm{T})}$.}
\begin{questions}
\item $r=\qty{26571}{km}$ et $T=\res{\qty{4,31e4}{s}}$, soit \res{\text{11~h~59~min}} : la moitié de la durée de rotation de la Terre.
\item $v=\dfrac{2\pi r}{T}=\res{\qty{3,87e3}{m.s^{-1}}}$.
\item \res{\text{Non}} : sa période vaut environ la moitié d'un jour sidéral ; il repasse au-dessus du même point
\res{\text{deux fois}} par jour sidéral.
\end{questions}
\end{corrige}
\end{blocLivre}
\begin{blocLivre}{cor-16-10}
\begin{corrige}{10}{Observer la Terre}
\demarche{Période par la troisième loi ; nombre de tours $=86400/T$.}
\begin{questions}
\item $r=\qty{7157}{km}$ : $T=\res{\qty{6,03e3}{s}}$ (100,5~min) et $\dfrac{86400}{6030}=\res{14{,}3}$ tours par jour.
\item Pendant que le satellite tourne autour des pôles, la Terre tourne sous lui : chaque passage survole une bande
différente, et toute la surface est couverte en quelques jours.
\item \res{\text{Non}} : un satellite géostationnaire doit être sur une orbite équatoriale, à \num{35800}~km d'altitude.
\end{questions}
\end{corrige}
\end{blocLivre}
\begin{blocLivre}{cor-16-11}
\begin{corrige}{11}{Lancer depuis Kourou}
\demarche{Comparer la vitesse « offerte » par la rotation terrestre à la vitesse de satellisation.}
\begin{questions}
\item $v_1=\sqrt{\dfrac{\num{3.98e14}}{\num{6.371e6}}}=\res{\qty{7,90e3}{m.s^{-1}}}$.
\item $\dfrac{2\pi R_\mathrm{T}}{T}=\qty{465}{m.s^{-1}}$ : à Kourou, $\res{\qty{463}{m.s^{-1}}}$ ; à Baïkonour, $\res{\qty{323}{m.s^{-1}}}$.
\item Soit \res{5{,}9~\%} de $v_1$ à Kourou et \res{4{,}1~\%} à Baïkonour : proche de l'équateur, Kourou offre le gain maximal,
d'où des économies de carburant.
\item La Terre tourne d'ouest en est : en lançant vers l'est, on additionne sa vitesse de rotation à celle du lanceur.
\end{questions}
\end{corrige}
\end{blocLivre}
\begin{blocLivre}{cor-16-12}
\begin{corrige}{12}{Le télescope spatial Hubble}
\demarche{$v=\sqrt{GM/r}$ : plus l'orbite est basse, plus la vitesse est grande et la période courte.}
\begin{questions}
\item Seule l'attraction terrestre (le freinage est très faible) : mouvement circulaire uniforme ;
$v=\res{\qty{7,59e3}{m.s^{-1}}}$ et $T=\res{95{,}4~\text{min}}$.
\item $\vec{v}$ tangent à l'orbite ; $\vec{a}$ dirigé vers le centre de la Terre, de norme $v^2/r=\qty{8,3}{m.s^{-2}}$.
\item À 500~km, $r$ diminue : $v$ \res{\text{augmente}} ($\qty{7,61e3}{m.s^{-1}}$) et $T$ \res{\text{diminue}} (94,5~min).
\item Le frottement fait perdre de l'énergie et abaisse l'orbite ; sur une orbite plus basse, la gravitation impose une
vitesse plus grande : le télescope « tombe » vers une orbite plus rapide.
\end{questions}
\schema[corrige]{rev-16-ex12-cor}
\legende{Orbite circulaire modélisée ; la vitesse est tangente et l’accélération centripète. Dessin sans échelle.}
\end{corrige}
\end{blocLivre}
\begin{blocLivre}{cor-16-13}
\begin{corrige}{13}{Les deux lunes de Mars}
\demarche{$M=\dfrac{4\pi^2r^3}{GT^2}$ pour chaque lune ; l'orbite synchrone a la période de rotation de Mars.}
\begin{questions}
\item Deuxième loi dans le repère de Frenet, puis $T=\dfrac{2\pi r}{v}$ : $\dfrac{T^2}{r^3}=\dfrac{4\pi^2}{G\,M_\mathrm{Mars}}$.
\item Phobos : $\res{\qty{6,42e23}{kg}}$ ; Deimos : $\res{\qty{6,43e23}{kg}}$. Les deux valeurs concordent : la troisième loi est
vérifiée et la masse de Mars est bien déterminée.
\item $r=\left(\dfrac{GM\,T^2}{4\pi^2}\right)^{1/3}$ avec $T=\qty{24,6}{h}$ : $r=\res{\qty{2,04e4}{km}}$.
\item Phobos est plus proche que l'orbite synchrone : il fait le tour de Mars (7,66~h) plus vite que Mars ne tourne sur
elle-même (24,6~h). Il dépasse donc la rotation de la planète et semble se déplacer d'ouest en est dans le ciel martien.
\end{questions}
\end{corrige}
\end{blocLivre}
\begin{blocLivre}{cor-16-14}
\begin{corrige}{14}{Les sept planètes de TRAPPIST-1}
\demarche{Convertir les périodes en années, puis calculer $T^2/a^3$ ; ce rapport vaut $M_\odot/M$.}
\begin{questions}
\item $T^2/a^3$ est le même pour tous les corps en orbite autour d'un astre ; il vaut $4\pi^2/(GM)$ et dépend donc de la
masse de l'astre central.
\item Pour b, e et h : $\dfrac{T^2}{a^3}\approx\res{11{,}1~\text{an}^2.\text{ua}^{-3}}$ dans les trois cas : la troisième loi est
vérifiée dans ce système.
\item $M=\dfrac{M_\odot}{11{,}1}=\res{0{,}090\,M_\odot}$ : une étoile de moins d'un dixième de la masse du Soleil.
\item TRAPPIST-1, naine rouge, est beaucoup moins lumineuse que le Soleil : il faut être beaucoup plus près d'elle pour
recevoir la même puissance par unité de surface.
\end{questions}
\end{corrige}
\end{blocLivre}
\begin{blocLivre}{cor-16-15}
\begin{corrige}{15}{L'orbite de transfert}
\demarche{Le demi-grand axe vaut $(r_1+r_2)/2$ ; la durée du transfert est une demi-période de l'ellipse.}
\begin{questions}
\item $v_1=\sqrt{\dfrac{GM_\mathrm{T}}{r_1}}=\res{\qty{7,72e3}{m.s^{-1}}}$ et $v_2=\sqrt{\dfrac{GM_\mathrm{T}}{r_2}}=\res{\qty{3,07e3}{m.s^{-1}}}$.
\item Le grand axe va du périgée à l'apogée : $2a=r_1+r_2$, soit $a=\res{\qty{2,44e4}{km}}$.
\item $T=2\pi\sqrt{\dfrac{a^3}{GM_\mathrm{T}}}=\qty{3,80e4}{s}$ ; le transfert dure une demi-période : \res{\qty{5,3}{h}}.
\item Le satellite s'éloigne de la Terre : d'après la loi des aires, il parcourt des arcs de plus en plus courts pendant
des durées égales ; sa vitesse diminue jusqu'à l'apogée.
\end{questions}
\schema[corrige]{16-ex15}
\end{corrige}
\end{blocLivre}
